\subsection{Erweiterte Koeffizientenmatrix}
\label{subsec:linear_systems_augmented_matrix}

Die Koeffizienten und Absolutglieder eines LGS lassen sich in einem rechteckigen
Zahlenschema anordnen. Die ersten $n$ Spalten bilden die
\emph{Koeffizientenmatrix} \EnglishTerm{coefficient matrix} $A$, die letzte
Spalte enthält den Vektor $b$ der Absolutglieder. Zusammen bilden sie die
\emph{erweiterte Koeffizientenmatrix} \EnglishTerm{augmented matrix}
$[A\mid b]$.

\begin{center}
  \CachedTikzFigure{content/resources/chapters/linear_systems/tikz/augmented_matrix_anatomy.tex}
\end{center}

Diese Notation verzichtet auf eine explizite Nennung der Unbekannten $x_j$.
Ihre Anzahl entspricht der Anzahl der Spalten von $A$.

Die Symbole $A$ und $b$ dienen hier ausschließlich als kompakte Notation;
Matrizen werden später ausführlich behandelt.

\begin{ExampleBox}{Vom LGS zur erweiterten Matrix}{linear_systems_augmented_matrix_example}
  Wir überführen die LGS in ihre erweiterte Koeffizientenmatrix $[A\mid b]$.
  \begin{itemize}[leftmargin=1.4em,itemsep=0.8em]
    \item
      \noindent
      \begin{minipage}[c]{0.47\linewidth}
        \[
          \begin{aligned}
             x_1-x_2       &=1,\\
            2x_1+3x_2+4x_3&=5
          \end{aligned}
        \]
      \end{minipage}%
      \begin{minipage}[c]{0.08\linewidth}
        \centering $\Longleftrightarrow$
      \end{minipage}%
      \begin{minipage}[c]{0.42\linewidth}
        \begin{ClozeArea}[12]{4}
          \centering
          \[
            \left[
              \begin{array}{rrr|r}
                1&-1&0&1\\
                2& 3&4&5
              \end{array}
            \right]
          \]
        \end{ClozeArea}
      \end{minipage}

    \item
      \noindent
      \begin{minipage}[c]{0.47\linewidth}
        \[
          \begin{aligned}
             2x_1+x_2&= 0,\\
              -x_1+3x_2&= 7,\\
                   4x_2&=-2
          \end{aligned}
        \]
      \end{minipage}%
      \begin{minipage}[c]{0.08\linewidth}
        \centering $\Longleftrightarrow$
      \end{minipage}%
      \begin{minipage}[c]{0.42\linewidth}
        \begin{ClozeArea}[12]{6}
          \centering
          \[
            \left[
              \begin{array}{rr|r}
                 2&1& 0\\
                -1&3& 7\\
                 0&4&-2
              \end{array}
            \right]
          \]
        \end{ClozeArea}
      \end{minipage}
  \end{itemize}
\end{ExampleBox}
